\section{Keeping the full phase correction}
\label{v6phase:section}

Let $c=C_{\mathrm E}$, $h=|H|$, $w=h^2/(2n)$ and $e=e^{-w}$.
Write the conserved identity as
\[
 R=U+n(1-e)\left(c\rho-\frac{\phi(0)D}{6}\right),\qquad
 U=(n+H^2)a+\frac{n\phi(0)D}{6}+n^{3/2}\mathcal L.
\]
Assume $a,D,\mathcal L\ge0$, $n\ge N>0$, and a valid positive lower
bound $b\le c\rho$. Since $U\ge n\phi(0)D/6$, this gives
\begin{equation}\label{v6phase:exact}
 U\le nb+e^{h^2/(2n)}(R-nb).
\end{equation}
Unlike the earlier constant phase price, this estimate keeps the dependence
of $R$ on the threshold. It requires no small-defect assumption.

\begin{lemma}[Only the floor and infinite-index limit are needed]
\label{v6phase:index}
For any real $r$, $b>0$, $h\ge0$ and $n\ge N>0$,
\begin{equation}\label{v6phase:endpoints}
 nb+e^{h^2/(2n)}(r-nb)
 \le\max\left\{r-\frac{bh^2}{2},\;
 Nb+e^{h^2/(2N)}(r-Nb)\right\}.
\end{equation}
This is an inequality for every real index $n\ge N$, hence for all integers.
\end{lemma}
\begin{proof}
The case $h=0$ is immediate. Put $a_0=h^2/2>0$, $x=1/n$, and extend
\[
 f(x)=re^{a_0x}-b\frac{e^{a_0x}-1}{x}
\]
continuously to $x=0$, with $f(0)=r-ba_0$. Its derivative has the sign of
\[
 r-ba_0\,j(a_0x),\qquad
 j(z)=\frac{z-1+e^{-z}}{z^2}
     =\int_0^1(1-s)e^{-zs}\,ds.
\]
Indeed $f'(x)=a_0e^{a_0x}[r-ba_0j(a_0x)]$.
The integral shows that $j$ strictly decreases on $[0,\infty)$.
Thus the derivative can change sign only from negative to positive.
There is no interior maximum on $[0,1/N]$; the two endpoint values give
\eqref{v6phase:endpoints}. No monotonicity of the full expression in $n$
is presumed.
\end{proof}

The infinite-index endpoint can be dispensed with when an upper bound
on a nonnegative resource is wanted. If $f(0)=r-ba_0>0$, then
$r>ba_0$ and $j\le1/2$ make $f'>0$, so $f(1/N)>f(0)$.
Otherwise $f(0)\le0$. Thus the right side of
\eqref{v6phase:endpoints} is at most the positive part of its finite
endpoint. Keeping that positive part is essential; the expression need
not be monotone in $n$ when its values are negative.

Suppose now a valid raw envelope is
\[
 R\le r(h)=A+Fh-qh^2\le T<Nb,\qquad F\ge0,\quad q>0.
\]
The first inequality in \eqref{v6phase:exact} is increasing in $R$, so
replace it by $r(h)$. At the finite endpoint, $r(h)-Nb<0$.
Consequently the lower Taylor bound $e^z\ge1+z+z^2/2$ gives an
\emph{upper} resource bound. Define
\begin{align*}
 P_\infty(h)&=A+Fh-(q+b/2)h^2,\\
 P_N(h)&=r(h)+(r(h)-Nb)
                    \left(\frac{h^2}{2N}+\frac{h^4}{8N^2}\right).
\end{align*}
Then, uniformly for all $n\ge N$ and all thresholds,
\begin{equation}\label{v6phase:polynomial}
 U\le\max\{0,P_N(h)\}.
\end{equation}
The infinite-index polynomial has exact maximum
$A+F^2/[4(q+b/2)]$ and remains a useful comparison.
The finite-index polynomial has degree six and negative leading coefficient.
On $r(h)<0$ the
nonnegative raw conserved quantity cannot occur, so a certified finite
interval containing $\{h\ge0:r(h)\ge0\}$ suffices.

There is a simple analytic test for a unique maximum; no positivity
subdivision is necessary. Put
\[
 z=\frac{F^2}{8q^2N}.
\]
If $z+z^2<1$, then $P_N$ is strictly concave on $[0,\infty)$.
Indeed, writing $v(h)=h^2/(2N)+h^4/(8N^2)$ gives
\[
 P_N''=-2q(1+v)+2(F-2qh)v'+(r-Nb)v''.
\]
Here $v''>0$ and $r-Nb<0$. For $h\ge F/(2q)$ every term is
nonpositive and the first is negative. Before that point,
$h^2/(2N)\le z$ and
$h(F-2qh)\le F^2/(8q)$, whence
\[
 P_N''\le-2q+\frac{F^2}{4qN}(1+z)
          =-2q(1-z-z^2)<0.
\]
Since $P_N'(0)=F\ge0$, its unique maximum is attained at $h=0$ when
$F=0$, and otherwise at its unique positive stationary point. Rational bisection of the
derivative followed by an outward monomial bound encloses that value.
The check of $z+z^2<1$ is a scalar inequality, not an assumption of
sampled concavity. The old 24-cell coefficients already have $z<0.325$;
decreasing $F$ or increasing $q$ preserves this gate.

For a formal coefficient interval containing $c$, take $A,F,q,b$ affine
in that coefficient, holding the same-law moment bounds fixed. Each of
the two displayed polynomials is then affine before maximization over
$h$. Their suprema, and their maximum, are convex in the coefficient;
its two endpoints again suffice. The strict gate $T<Nb$ must hold for
the chosen envelopes before reversing the exponential estimate.

This is a stronger way to bound $U$, not a new raw-remainder estimate.
The raw ceiling $T$ still supplies any earlier density license. Neither
the low/high Fourier estimates nor the finite base are removed by this
calculation.
